% =============================================================================
% 02_einleitung.tex — Einleitung
% =============================================================================
% Die Einleitung führt die Lesenden in die Arbeit ein. Sie enthält:
%   1. Grundlagen und Definition des Themas
%   2. Ziele der Arbeit
%   3. Technisches Vorgehen / Anforderungsspezifikation
%   4. Aufbau der Arbeit
% =============================================================================

\chapter{Einleitung}

% --- 1. Grundlagen und Definition des Themas ---------------------------------
% Erkläre den fachlichen Kontext. Was ist ein Geotracker?
% Welche technologischen Grundlagen (GPS, IoT, Cloud, Web) sind relevant?
% Definiere zentrale Begriffe, die in der Arbeit verwendet werden.
\section{Grundlagen und Themendefinition}

<<Hier den technischen und fachlichen Hintergrund des Geotrackers beschreiben.
Zentrale Begriffe definieren (z.\,B.\ GPS, GNSS, MQTT, REST-API, Progressive
Web App).>>

% --- 2. Ziele der Arbeit -----------------------------------------------------
% Formuliere klare, messbare Ziele (z. B. als Forschungsfragen oder Muss-Ziele).
\section{Ziele der Arbeit}

<<Hier die konkreten Ziele und Forschungsfragen der IDPA formulieren.
Beispiel: Kann ein kostengünstiger GPS-Tracker mit Echtzeit-Webdarstellung
realisiert werden?>>

% --- 3. Technisches Vorgehen / Anforderungsspezifikation ---------------------
% Beschreibe den gewählten Entwicklungsansatz (z. B. agil, iterativ).
% Liste funktionale und nicht-funktionale Anforderungen auf.
\section{Technisches Vorgehen und Anforderungsspezifikation}

\subsection{Funktionale Anforderungen}
% Was muss das System können? (Use-Cases, User Stories o. ä.)
<<Funktionale Anforderungen als Liste oder Tabelle einfügen.>>

\subsection{Nicht-funktionale Anforderungen}
% Leistung, Sicherheit, Zuverlässigkeit, Wartbarkeit usw.
<<Nicht-funktionale Anforderungen beschreiben.>>

% --- 4. Aufbau der Arbeit ----------------------------------------------------
% Erkläre kurz, was in den folgenden Kapiteln behandelt wird.
\section{Aufbau der Arbeit}

Die vorliegende Arbeit gliedert sich wie folgt:
\begin{description}
  \item[Kapitel~\ref{chap:hauptteil} – Hauptteil]
    Beschreibt die Entwicklung in den drei technischen Domänen Firmware,
    Cloud-Konfiguration und Web-Frontend.
  \item[Kapitel~\ref{chap:fazit} – Fazit]
    Bewertet die erzielten Ergebnisse kritisch und beantwortet die
    Forschungsfragen.
  \item[Anhang]
    Enthält ergänzende Unterlagen, detaillierte Auswertungen und
    unveröffentlichte Quellen.
\end{description}
